\documentclass[11pt]{article}
\usepackage[T1]{fontenc}
\usepackage[utf8]{inputenc}
\usepackage{lmodern}
\usepackage{amsmath,amssymb,amsthm,mathtools}
\usepackage[margin=1.1in]{geometry}
\usepackage{microtype}
\usepackage{enumitem}
\usepackage{hyperref}
\usepackage{xcolor}
\hypersetup{colorlinks=true,linkcolor=blue!50!black,citecolor=blue!50!black,urlcolor=blue!50!black}

\newtheoremstyle{plainstyle}{1em}{1em}{\normalfont}{} {\scshape}{.}{0.5em}{}
\newtheoremstyle{defstyle}{1em}{1em}{\normalfont}{} {\itshape}{.}{0.5em}{}
\theoremstyle{plainstyle}
\newtheorem{theorem}{Theorem}[section]
\newtheorem{proposition}[theorem]{Proposition}
\newtheorem{corollary}[theorem]{Corollary}
\theoremstyle{defstyle}
\newtheorem{definition}[theorem]{Definition}
\newtheorem{remark}[theorem]{Remark}
\newtheorem{summary}[theorem]{Summary}

\newcommand{\A}{\mathcal{A}}
\newcommand{\C}{\mathbb{C}}
\newcommand{\Htau}{\mathcal{H}_{\tau}}
\newcommand{\Dtau}{\mathcal{D}_{\tau}}
\newcommand{\Omegatau}{\Omega_{\tau}}

\title{Yang--Mills Boundary Packages with Instanton Completion and Screening\\
\large Internal screening domination as the mass-gap reduction target}
\author{}
\date{March 2026}

\begin{document}
\maketitle

\begin{summary}[Executive synthesis]
The instanton-completed boundary package of a holomorphic--topological Yang--Mills theory
produces, after passing to the Koszul-dual defect algebra, a finite family of central screening classes.
Once a positive vacuum state is chosen, these screenings define a canonical positive quadratic form on the
GNS Hilbert space of the defect algebra.  The key point is that this quadratic form is \emph{internal} to the
formalism: it depends only on the defect algebra, the vacuum state, and the screening classes.
The untwisted mass-gap problem is therefore reduced to a single analytic estimate: prove that the physical
Yang--Mills Hamiltonian dominates this internal screening form on the algebraic core.
Under this domination estimate, a positive mass gap follows by an elementary min--max argument.
\end{summary}

\section{From instanton completion to an internal screening Hamiltonian}

The bar--cobar programme for boundary conditions in holomorphic--topological gauge theory
produces an instanton-completed Koszul-dual defect algebra $\A_{\mathrm{def}}$, together with a visible center
obtained by quotienting by a central screening ideal.
For the mass-gap problem, the point is not merely that the visible center controls first-order couplings,
but that the screening classes themselves already define a canonical positive operator on a Hilbert completion
of the defect algebra.
This operator will be the internal comparison object for the untwisted Hamiltonian.

\begin{definition}[Screening state datum]
Let $\A_{\mathrm{def}}$ be a unital $*$-algebra and let
\[
\tau\colon \A_{\mathrm{def}}\longrightarrow \C
\]
be a positive normalized state.  Let
\[
\mathbf s=(s_1,\dots,s_r)\subset Z(\A_{\mathrm{def}})
\]
be a finite family of degree-zero central screening classes.  We say that
$(\A_{\mathrm{def}},\tau,\mathbf s)$ is a \emph{screening state datum} if:
\begin{enumerate}[label=(\roman*)]
\item each $s_i$ is \emph{$\tau$-bounded}, i.e. there exists $C_i\ge 0$ such that
\[
\tau\bigl((a s_i)^*(a s_i)\bigr)\le C_i^2\,\tau(a^*a)
\qquad\text{for all }a\in \A_{\mathrm{def}};
\]
\item each screening annihilates the vacuum in the GNS representation, equivalently
\[
\tau(s_i^*s_i)=0
\qquad\text{for all }i.
\]
\end{enumerate}
\end{definition}

Let $(\pi_\tau,\Htau,\Omegatau)$ be the GNS representation of $\tau$, and let
\[
\Dtau:=\pi_\tau(\A_{\mathrm{def}})\Omegatau
\subset \Htau
\]
be the algebraic core.
For each $i$, define an operator on $\Dtau$ by right multiplication,
\[
S_i[a]:=[a s_i],
\qquad [a]:=\pi_\tau(a)\Omegatau.
\]
By $\tau$-boundedness, $S_i$ extends uniquely to a bounded operator on $\Htau$.
Because each $s_i$ is central, $S_i$ commutes with the left action of $\A_{\mathrm{def}}$.
We set
\[
H_{\mathrm{scr}}^{\mathrm{int}}:=\sum_{i=1}^r S_i^*S_i,
\]
and call it the \emph{internal screening Hamiltonian}.
We further define the associated screening quadratic form on algebraic vectors by
\[
\mathfrak s_{\mathbf s}(a):=\sum_{i=1}^r \tau\bigl((a s_i)^*(a s_i)\bigr).
\]

\begin{proposition}[Internal screening form]
For every screening state datum:
\begin{enumerate}[label=(\alph*)]
\item the operators $S_i$ are bounded and pairwise commuting;
\item $H_{\mathrm{scr}}^{\mathrm{int}}$ is bounded and nonnegative;
\item $H_{\mathrm{scr}}^{\mathrm{int}}\Omegatau=0$;
\item for every $a\in \A_{\mathrm{def}}$,
\[
\mathfrak s_{\mathbf s}(a)=\bigl\langle [a],H_{\mathrm{scr}}^{\mathrm{int}}[a]\bigr\rangle_{\Htau}.
\]
\end{enumerate}
\end{proposition}

\begin{proof}
The only nontrivial point is boundedness, which is exactly the $\tau$-boundedness hypothesis.
Centrality gives commutativity.  Positivity of $H_{\mathrm{scr}}^{\mathrm{int}}$ is immediate from the formula
$\sum_i S_i^*S_i$.  Since
\[
\|S_i\Omegatau\|^2=\tau(s_i^*s_i)=0,
\]
we have $S_i\Omegatau=0$ and hence $H_{\mathrm{scr}}^{\mathrm{int}}\Omegatau=0$.
Finally,
\[
\bigl\langle [a],S_i^*S_i[a]\bigr\rangle=\|S_i[a]\|^2
=\|[a s_i]\|^2=\tau\bigl((a s_i)^*(a s_i)\bigr),
\]
and summing over $i$ gives the stated identity.
\end{proof}

\section{The screening gap and its algebraic form}

The screening Hamiltonian controls the directions in Hilbert space that remain visible after quotienting by the screening ideal.
Its spectral positivity can be detected purely on algebraic vectors.
Let $P_0$ denote the orthogonal projection onto the vacuum line $\C\Omegatau$.

\begin{proposition}[Algebraic and operator formulations of the screening gap]
The following are equivalent for a fixed constant $\mu>0$:
\begin{enumerate}[label=(\roman*)]
\item the operator inequality
\[
H_{\mathrm{scr}}^{\mathrm{int}}\ge \mu^2(1-P_0)
\]
holds on $\Htau$;
\item the algebraic inequality
\[
\mathfrak s_{\mathbf s}(a)\ge \mu^2\,\bigl\|[a]-\tau(a)\Omegatau\bigr\|^2
\qquad\text{for all }a\in \A_{\mathrm{def}}
\]
holds on the dense core $\Dtau$.
\end{enumerate}
\end{proposition}

\begin{proof}
The vectors $[a]-\tau(a)\Omegatau$ span a dense subspace of the orthogonal complement of the vacuum.
Restricting the operator inequality to this dense subspace gives the algebraic inequality.
Conversely, the algebraic inequality is a lower bound for the quadratic form of the bounded operator
$H_{\mathrm{scr}}^{\mathrm{int}}$ on a form core, so it extends to all of $\Htau$ by continuity.
\end{proof}

\begin{definition}[Internal screening gap]
We say that $(\A_{\mathrm{def}},\tau,\mathbf s)$ has \emph{screening gap} at least $\mu>0$
if the equivalent conditions above hold.
\end{definition}

\begin{remark}[Why this is internal]
Every quantity in the algebraic inequality
\[
\mathfrak s_{\mathbf s}(a)\ge \mu^2\,\bigl\|[a]-\tau(a)\Omegatau\bigr\|^2
\]
is defined using only the defect algebra $\A_{\mathrm{def}}$, the vacuum state $\tau$, and the central screening sequence $\mathbf s$.
Thus the screening gap is a property of the instanton-completed boundary package itself.
No external spacetime Hamiltonian has entered yet.
\end{remark}

\section{The screening domination estimate and the mass gap}

We now formulate the one remaining analytic estimate.
Let $H_{\mathrm{YM}}$ be a nonnegative self-adjoint operator on $\Htau$ with
\[
H_{\mathrm{YM}}\Omegatau=0,
\qquad
\Dtau\subset \operatorname{Dom}(H_{\mathrm{YM}}^{1/2}).
\]
The operator $H_{\mathrm{YM}}$ should be thought of as the untwisted Yang--Mills Hamiltonian once an untwisting bridge has been constructed.

\begin{definition}[Algebraic screening domination estimate]
We say that $H_{\mathrm{YM}}$ satisfies the \emph{algebraic screening domination estimate}
with constants $(\alpha,\varepsilon)$ if
\[
\bigl\langle [a],H_{\mathrm{YM}}[a]\bigr\rangle
\ge
\alpha\,\mathfrak s_{\mathbf s}(a)
-\varepsilon\,\bigl\|[a]-\tau(a)\Omegatau\bigr\|^2
\qquad\text{for all }a\in \A_{\mathrm{def}}.
\]
\end{definition}

\begin{theorem}[Mass-gap reduction theorem]
Assume that $(\A_{\mathrm{def}},\tau,\mathbf s)$ has screening gap at least $\mu>0$.
Assume further that $H_{\mathrm{YM}}$ satisfies the algebraic screening domination estimate
with constants $\alpha>0$ and $0\le \varepsilon<\alpha\mu^2$.
Then
\[
\operatorname{Spec}(H_{\mathrm{YM}})\cap (0,\alpha\mu^2-\varepsilon)=\varnothing.
\]
In particular, $H_{\mathrm{YM}}$ has a positive mass gap at least $\alpha\mu^2-\varepsilon$.
\end{theorem}

\begin{proof}
Let $\psi=[a]-\tau(a)\Omegatau\in \Dtau\cap (\C\Omegatau)^{\perp}$.
Since $H_{\mathrm{YM}}\Omegatau=0$, we have
\[
\langle \psi,H_{\mathrm{YM}}\psi\rangle
=
\langle [a],H_{\mathrm{YM}}[a]\rangle.
\]
Applying the domination estimate and then the screening-gap inequality gives
\[
\langle \psi,H_{\mathrm{YM}}\psi\rangle
\ge
\alpha\,\mathfrak s_{\mathbf s}(a)-\varepsilon\,\|\psi\|^2
\ge
(\alpha\mu^2-\varepsilon)\,\|\psi\|^2.
\]
Thus the quadratic form of $H_{\mathrm{YM}}$ is bounded below by $\alpha\mu^2-\varepsilon$ on a dense core for the orthogonal complement of the vacuum.
The min--max principle excludes spectrum in $(0,\alpha\mu^2-\varepsilon)$.
\end{proof}

\begin{corollary}[Exact reduction target]
Assume a strictly positive screening gap has been established.
Then the untwisted Yang--Mills mass-gap problem is reduced to proving the single estimate
\[
\bigl\langle [a],H_{\mathrm{YM}}[a]\bigr\rangle
\ge
\alpha\sum_{i=1}^r \tau\bigl((a s_i)^*(a s_i)\bigr)
-\varepsilon\,\bigl\|[a]-\tau(a)\Omegatau\bigr\|^2
\qquad\text{for all }a\in \A_{\mathrm{def}},
\]
with some $\alpha>0$ and $0\le \varepsilon<\alpha\mu^2$.
Everything on the right-hand side is internal to the instanton-completed boundary package.
\end{corollary}

\begin{remark}[Frontier]
The reduction theorem cleanly separates the problem into two parts.
The first is algebraic: construct the defect algebra, the screening sequence, and a positive vacuum state with strictly positive screening gap.
The second is analytic: construct the untwisting bridge and prove domination of the physical Hamiltonian by the internal screening form.
The formalism now provides an explicit and verifiable target for the analytic step.
\end{remark}

\end{document}
